\chapter*{Course Map: Learn One Branch at a Time}
\addcontentsline{toc}{chapter}{Course Map: Learn One Branch at a Time}
\markboth{Course Map}{ }

Before the detailed Machine-Learning Map, use this one-page route to answer only one question: \emph{what should I learn next?} The full field taxonomy appears after Chapter~1, once the basic vocabulary has been introduced.

\begin{mlfigure}
\begin{tikzpicture}[
  every node/.append style={execute at begin node={\hyphenpenalty=10000\relax}},
  stage/.style={draw=MLBlue,fill=MLPaleBlue,rounded corners=2mm,text width=12.75cm,minimum height=1.38cm,align=left,inner xsep=4mm,inner ysep=2.4mm,font=\scriptsize\color{MLNavy}},
  gate/.style={draw=MLNavy,fill=MLNavy,text=white,rounded corners=2mm,text width=6.8cm,minimum height=0.78cm,align=center,font=\small\bfseries},
  branch/.style={rounded corners=2mm,text width=5.90cm,minimum height=2.85cm,align=left,inner xsep=3.4mm,inner ysep=2.5mm,font=\scriptsize\color{MLNavy}},
  sup/.style={branch,draw=MLBlue,fill=MLPaleBlue},
  unsup/.style={branch,draw=MLTeal,fill=MLPaleTeal},
  other/.style={draw=MLOrange,fill=MLPaleOrange,rounded corners=2mm,text width=12.75cm,minimum height=1.45cm,align=left,inner xsep=4mm,inner ysep=2.4mm,font=\scriptsize\color{MLNavy}},
  finish/.style={draw=MLGreen,fill=MLPaleTeal,rounded corners=2mm,text width=12.75cm,minimum height=1.30cm,align=left,inner xsep=4mm,inner ysep=2.4mm,font=\scriptsize\color{MLNavy}},
  branchlabel/.style={fill=white,inner xsep=2pt,inner ysep=1pt,font=\tiny\bfseries\color{MLMuted}},
  arr/.style={-{Stealth[length=2mm]},thick,draw=MLTeal},
  flow/.style={thick,draw=MLTeal,line cap=round,line join=round},
  junction/.style={circle,fill=MLTeal,inner sep=1.2pt}
]
\node[stage] (a) at (0,9.15) {{\normalsize\bfseries\color{MLBlue}Stage 1: Orientation and tools}\hfill{\footnotesize\bfseries Chapters 1--2}\par
  \smallskip\textbf{Learn:} problem vocabulary \textbullet\ the end-to-end workflow \textbullet\ Python and project organization.\par
  \textbf{Ready when:} you can state the problem clearly and run the course environment.};
\node[stage] (b) at (0,6.95) {{\normalsize\bfseries\color{MLBlue}Stage 2: Data and shared foundations}\hfill{\footnotesize\bfseries Chapters 3--5}\par
  \smallskip\textbf{Learn:} data quality, safe splitting and preprocessing, statistics, probability, geometry, scaling, and optimization.\par
  \textbf{Ready when:} the data is trustworthy and you can explain the mathematical ideas reused by later models.};
\node[gate] (g) at (0,5.10) {Which learning setting fits the task?};

\node[sup] (s1) at (-3.45,2.45) {{\centering\bfseries\color{MLBlue}Stage 3A: Supervised workflow\par}{\centering\bfseries Chapters 6--8\par}\smallskip
  \textbf{Learn:} baselines, metrics, honest validation, and a sealed final test.\par
  \textbf{Outcome:} a defensible selection procedure.};
\node[sup] (s2) at (-3.45,-0.70) {{\centering\bfseries\color{MLBlue}Stage 3B: Supervised model families\par}{\centering\bfseries Chapters 9--14\par}\smallskip
  \textbf{Compare:} linear/logistic; KNN/naive Bayes; trees/ensembles/SVMs.\par
  \textbf{Outcome:} the simplest adequate model.};

\node[unsup] (u1) at (3.45,2.45) {{\centering\bfseries\color{MLTeal}Stage 4A: Unsupervised workflow\par}{\centering\bfseries Chapter 15\par}\smallskip
  \textbf{Learn:} evaluation through stability, geometry, reconstruction, and domain evidence.\par
  \textbf{Outcome:} a defensible evaluation plan.};
\node[unsup] (u2) at (3.45,-0.70) {{\centering\bfseries\color{MLTeal}Stage 4B: Unsupervised model families\par}{\centering\bfseries Chapters 16--18\par}\smallskip
  \textbf{Compare:} clustering, representation learning, and anomaly detection.\par
  \textbf{Outcome:} structure supported by evidence, not assumed to be truth.};

\node[other] (c) at (0,-3.70) {{\normalsize\bfseries\color{MLOrange}Stage 5: Advanced Concepts}\hfill{\footnotesize\bfseries Chapters 19--20}\par
  \smallskip\textbf{Learn:} semi-supervised and active learning for scarce labels \textbullet\ reinforcement learning for actions and rewards.\par
  \textbf{Ready when:} you can explain how the workflow changes with the available feedback.};
\node[finish] (d) at (0,-5.95) {{\normalsize\bfseries\color{MLGreen}Capstone and Epilogue: integrate, defend, and extend}\par
  \smallskip\textbf{Integrate:} choose the branch \textbullet\ build a leakage-safe pipeline \textbullet\ evaluate fairly \textbullet\ communicate limitations.\par
  \textbf{Extend:} identify the next field direction that matches your goals.};

\draw[arr] (a.south) -- (b.north);
\draw[arr] (b.south) -- (g.north);
\coordinate (split) at ($(g.south)+(0,-0.40)$);
\draw[flow] (g.south) -- (split);
\node[junction] at (split) {};
\draw[arr] (split) -| node[pos=0.28,above,branchlabel] {Dataset is labeled} (s1.north);
\draw[arr] (split) -| node[pos=0.28,above,branchlabel] {Dataset is unlabeled} (u1.north);
\draw[arr] (s1.south) -- (s2.north);
\draw[arr] (u1.south) -- (u2.north);
\coordinate (merge) at ($(c.north)+(0,0.25)$);
\draw[flow] (s2.south) |- (merge);
\draw[flow] (u2.south) |- (merge);
\node[junction] at (merge) {};
\draw[arr] (merge) -- (c.north);
\draw[arr] (c.south) -- (d.north);
\end{tikzpicture}
\end{mlfigure}

\begin{keyidea}
Do not memorize the map. Use it as a routing rule. First learn the workflow for a learning signal; only then compare model families within that workflow. The capstone asks you to choose and defend the route yourself.
\end{keyidea}